\documentclass[11pt,a4paper]{article}
\usepackage[T1]{fontenc}
\usepackage[utf8]{inputenc}
\usepackage[english]{babel}
\usepackage{amsmath,amssymb}
\usepackage[margin=2cm]{geometry}
\setlength{\parindent}{0pt}
\setlength{\emergencystretch}{3em}
\allowdisplaybreaks

\title{About \texttt{book\_notation.tex}}
\author{}
\date{}
\begin{document}
\maketitle
This note explains what the test data in \texttt{book\_notation.tex} is used for. That file is input for the engine in \texttt{main.c}; it is read by the test \texttt{test\_set\_theory.c} in this folder.

\begin{enumerate}
\item \textbf{Purpose.} The file holds 26 true statements written in the notation of the books: images $f(X)$ and preimages $f^{-1}(Y)$ of parts, relations given as sets of pairs $\{(x, y) \mid c\}$, their inverses and compositions, the diagonal $\{(x, x) \mid x \in S\}$, families of parts $X_{i}$ with $\bigcup$ and $\bigcap$, variables that range over a product $S \times T$, maps with values in a product, three examples, and (the last three) a relation and a map between the parts of a set given by a rule of set algebra. Each must be proved. They are not copies of the exercises of \texttt{maps.tex} or \texttt{correspondences.tex}: they are other facts of the same kinds, so that the reader and the unfolding of paragraphs 17p, 17q and 17r of \texttt{parser.c} are tested on statements the books do not contain.

\item \textbf{What is being tested.} That the notation is read and unfolded into conditions on elements, that the unfolded statement agrees with the one as written on finite models (paragraph 28h of \texttt{main.c}), and that chapter 28 proves it, with a certificate that \texttt{checker.c} accepts. The last three are unfolded into statements of chapter 14 (paragraph 17r), checked against the statement as written on small universes, and decided by regions (chapter 15), without a certificate because they have quantifiers over sets.

\item \textbf{How to run it.} From the top folder of the repository:
\begin{verbatim}
cc -std=c11 -O2 -Wall -Wextra -o /tmp/test_set_theory \
    tests/set_theory/test_set_theory.c && /tmp/test_set_theory
\end{verbatim}

\item \textbf{Expected result.} One \texttt{PASS} line for the file: every statement is proved, with its written proof. The false statements of the same kinds are in \texttt{false\_book\_relations.tex} and \texttt{false\_book\_maps.tex}.
\end{enumerate}
\end{document}
